\documentclass[../Main.tex]{subfiles}

\begin{document}

La pureza $\mu$ dada en la definición \eqref{def:pureza} es una medida es sencilla de calcular conociendo la forma exacta del operador de estado; sin embargo, nos gustaría una medida que dependa de alguna cantidad física medible. Con esta motivación en mente, consideremos un ensamble de partículas de espín $\frac{1}{2}$ y espín promedio definido (cantidad física medible); es decir, vamos a considerar un ensamble con dos niveles de energía.\\ Hacemos esta consideración, pues el desarrollo de esta tesis se basa en qubits que es un sistema de dos niveles de energía.

Dado a \eqref{def:Operador_densidad} definimos el espín promedio del ensamble como sigue:


\begin{equation}
    \bm{S}=\Tr(\rho \bm{\sigma}).
\end{equation}


Donde $\bm{\sigma}=(\sigma_x,\sigma_y,\sigma_z)$ son las matrices de Pauli.\\ 

Definimos una cantidad que mida de manera cuantitativa el desorden del sistema, la cual dependa del tipo de ensamble. Exigimos que dicha cantidad sea mínima (cero) para un ensamble puro y máxima para un ensamble completamente aleatorio.  

En termodinámica, la magnitud que cumple con estas propiedades es la \textit{entropía}. Motivados por esto, definimos la \textit{entropía en mecánica cuántica estadística} como \cite{Sakurai2011}
\begin{equation}
    S_{entropia} = -k_B \Tr(\rho \ln \rho),
    \label{def:entropia}
\end{equation}
donde $k_B$ es la constante de Boltzmann y $\rho$ es el operador densidad del sistema.\\
A continuación, buscaremos el operador densidad $\rho$ que describe el equilibrio térmico de un sistema con espín definido.

Consideremos la acción $F$:
\begin{align}
    F
    &=\Tr(S_{entropia}+\alpha_x S_x+\alpha_y S_y+\alpha_z S_z+\lambda \rho)\\
    &=\Tr(-\rho\ln{\rho}+\bm{\alpha}\cdot\Bar{S}+\lambda\rho).
\end{align}
Donde $S_i$ es el valor promedio de la i-esima componente del espín, es decir, $S_i=\Tr(\rho \sigma_i)$.\\
Sabemos que la entropía debe maximizarse en un estado de equilibrio, y dado que la entropía es cóncava en $\rho$ basta con obtener el punto donde $\delta F = 0$ para encontrar el estado $\rho$ que maximiza la entropía.
\begin{align}
0=\delta F
&= \Tr(\ln{\rho}+\Id+\bm{\alpha}\cdot \bm{\sigma}+\lambda\Id)\delta\rho.
\label{eq:multiplicadores}
\end{align}
entonces al ser para cualquier variación, necesariamente se requiere que:
\begin{equation}
    \rho = \frac{\e^{-\bm{\alpha}\cdot \bm{\sigma}}}{\e^{(1+\lambda)\Id}}.
    \label{eq:rho_wo_norm}
\end{equation}
Usando la restricción $\Tr(\rho)=1$, se obtiene
\begin{equation}
    \e^{(1+\lambda)\Id}=\Tr(\e^{-\bm{\alpha}\cdot\bm{\sigma}}).
    \label{eq:rho's_normaliz}
\end{equation}
Sustituyendo ~\eqref{eq:rho's_normaliz} en ~\eqref{eq:rho_wo_norm} se obtiene el operador densidad de un sistema con espín $\bm{S}$ y en equilibrio térmico.
\begin{equation}
    \rho = \frac{\e^{-\bm{\alpha}\cdot\bm{\sigma}}}{\Tr(\e^{-\bm{\alpha}\cdot\bm{\sigma}})}.
    \label{eq:rho_semifinal}
\end{equation}
Solo es necesario ahora incluir la restricción de espín para despejar los valores $\alpha_i$. Primero, recordando que:
\begin{equation}
(\bm{\sigma}\cdot\bm{\alpha})^{m}
=
\begin{cases}
|\bm{\alpha}|^{m} \;\Id,
& m \ \text{par}, \\[4pt]
|\bm{\alpha}|^{m}\,
\bm{\sigma}\cdot\hat{\bm{n}},
& m \ \text{impar},
\end{cases}
\end{equation}
Con ello, expandimos en serie de Taylor la exponencial:
\begin{align}
\e^{-\bm{\alpha}\cdot\bm{\sigma}}
&=
\sum_{m=0}^{\infty}
\frac{(-1)^m}{m!}
(\bm{\alpha}\cdot\bm{\sigma})^{m}
\notag \\
&=
\sum_{k=0}^{\infty}
\frac{|\bm{\alpha}|^{2k}}{(2k)!}\,\Id
-
\sum_{k=0}^{\infty}
\frac{|\bm{\alpha}|^{2k+1}}{(2k+1)!}
(\bm{\sigma}\cdot\hat{\bm{n}})
\notag \\
&=
\cosh|\bm{\alpha}|\,\Id
-
\frac{\sinh|\bm{\alpha}|}{|\bm{\alpha}|}
(\bm{\alpha}\cdot\bm{\sigma}).
\label{eq:expansion_exponencial}
\end{align}
Con este resultado, aplicamos la traza a la exponencial:
\begin{align}
\Tr\!\qty(\e^{-\bm{\alpha}\cdot\bm{\sigma}})
&=
\Tr\!\qty(
\cosh|\bm{\alpha}|\,\Id
-
\frac{\sinh|\bm{\alpha}|}{|\bm{\alpha}|}
(\bm{\alpha}\cdot\bm{\sigma})
)
\notag \\
&=
\cosh|\bm{\alpha}|\,\Tr(\Id)
-
\frac{\sinh|\bm{\alpha}|}{|\bm{\alpha}|}
\Tr(\bm{\alpha}\cdot\bm{\sigma})
\notag \\
&=
2\cosh|\bm{\alpha}|
-
\frac{\sinh|\bm{\alpha}|}{|\bm{\alpha}|}
\sum_i \alpha_i\; \cancelto{0}{\Tr(\sigma_i)}
\notag \\
&=
2\cosh|\bm{\alpha}|.
\label{eq:normalizacion_traza}
\end{align}
Sustituyendo ~\eqref{eq:normalizacion_traza} en ~\eqref{eq:rho_semifinal} y reacomodando:
\begin{equation}
    \rho = \frac{1}{2}\qty(\Id - \frac{\tanh{|\bm{\alpha}|}}{|\bm{\alpha}|})\bm{\alpha}\cdot\bm{\sigma}
    \label{eq:rho_semi_semifinal}
\end{equation}
Ahora, observando que
\begin{align}
    S_k 
    &= \Tr(\rho\,\sigma_k)
    = \Tr\!\qty[\frac{1}{2}\qty(\Id - \frac{\tanh |\bm{\alpha}|}{|\bm{\alpha}|}\sum_i \alpha_i\,\sigma_i )\sigma_k]
    \notag \\
    &= \frac{1}{2}\qty(\cancelto{0}{\Tr(\sigma_k)}
    - \frac{\tanh |\bm{\alpha}|}{|\bm{\alpha}|}
    \sum_i \alpha_i \Tr(\cancelto{2\delta_{ik}}{\sigma_i\sigma_k}\hspace{.9cm}))
    &= -\frac{\tanh |\bm{\alpha}|}{|\bm{\alpha}|}\,\alpha_k ,
\end{align}
y sustituyendo este resultado en la ecuación ~\eqref{eq:rho_semi_semifinal}, se obtiene finalmente
\begin{equation}
\rho = \frac{1}{2}\qty(\Id + \bm{S}\cdot \bm{\sigma}).
\label{eq:rho_final}
\end{equation}
Ahora, calculemos $\rho^2$ para hallar una relación con $\bm{S}$.
\begin{align*}
    \rho^2&=
    \frac{1}{4}\qty(\Id + \bm{S}\cdot \bm{\sigma})^2
    \notag \\
    &=\frac{1}{4}\qty(\Id+2\bm{S}\cdot\bm{\sigma}+|\bm{S}|^2\,\Id).
\end{align*}
Calculando la traza de $\rho^2$:
\begin{align}
    \Tr(\rho^2)&=
    \Tr[\frac{1}{4}\qty(\Id+2\,\bm{S}\cdot\bm{\sigma}+\bm{|S|}^2\,\Id)]
    \notag \\
    &= \frac{1}{4}\qty[\cancelto{2}{\Tr(\Id)}+2\sum_iS_i\;\cancelto{0}{\Tr(\sigma_i)}\;+2|\bm{S}|^2]
    \notag \\
    &=\frac{1}{2}\qty(1+|\bm{S}|^2).
\end{align}

De la última expresión se concluye que:

\begin{enumerate}
    \item Un estado puro corresponde a un vector de Bloch localizado en la superficie de la esfera de Bloch, es decir, $|\bm{S}|=1$.
    \item Un estado mixto corresponde a un vector de Bloch ubicado en el interior de la esfera, con $|\bm{S}|<1$.
    \item El estado completamente mezclado corresponde al caso $|\bm{S}|=0$, es decir, al origen de la esfera de Bloch.
\end{enumerate}

De esta manera, el espacio de operadores densidad es el conjunto de todas las $\rho$ de la forma de \eqref{eq:rho_final} y cada estado es identificado por su espín promedio $\bm{S}$ del sistema.
\begin{equation}
    \rho_{space}\equiv\qty{\rho=\frac{1}{2}\qty(\Id + \bm{S}\cdot \bm{\sigma})\;|\;\bm{S}\in \R^3 }
\end{equation}

\end{document}
